\section{ماژول mgie\_svd (MGIE-SVD Compressor)}
\label{sec:module_mgie_svd}

\subsection{معرفی و انگیزه}
در تقریب رتبه پایین وزن‌های ترنسفورمر به وسیله تجزیه مقادیر منفرد (\lr{SVD})، انتخاب رتبه بریده‌شده $r$ یکی از چالش‌های بنیادین است. در رویکردهای سنتی، یک رتبه ثابت و یکسان برای تمامی لایه‌ها اعمال می‌شود؛ این در حالی است که ماهیت ریاضی و پراکندگی مقادیر منفرد در لایه‌های مختلف به شدت ناهمگن است. برخی لایه‌ها دارای طیف مقادیر منفرد بسیار چوله‌دار با زوال سریع هستند (که با رتبه‌های بسیار پایین با دقت بالا قابل تقریب‌اند)، در حالی که لایه‌های دیگر دارای توزیع یکنواخت انرژی بوده و نیازمند حفظ ابعاد بالاتری هستند.

ماژول \lr{MGIE-SVD (Maximum Information Entropy SVD)} بر پایه اصل بیشینه‌سازی آنتروپی اطلاعاتی مقادیر منفرد طراحی شده است. این روش با محاسبه آنتروپی شانون برای توزیع احتمالاتی مقادیر منفرد در هر لایه، میزان پراکندگی طیفی لایه را به صورت خودکار ارزیابی نموده و رتبه تقریب را متناسب با آنتروپی اطلاعاتی آن لایه تخصیص می‌دهد. این امر موجب می‌شود لایه‌هایی که محتوای اطلاعاتی پیچیده‌تر و طیف یکنواخت‌تری دارند رتبه بالاتری دریافت کنند و لایه‌های با مقادیر منفرد غالب با رتبه‌های فشرده‌تر تقریب زده شوند.

\subsection{مبانی نظری و فرمول‌بندی ریاضی}
برای ماتریس وزن $\mathbf{W} \in \mathbb{R}^{d_{\text{out}} \times d_{\text{in}}}$ با ابعاد بیشینه $r_{\text{max}} = \min(d_{\text{out}}, d_{\text{in}})$، تجزیه مقادیر منفرد به صورت زیر انجام می‌شود:
\begin{equation}
\mathbf{W} = \mathbf{U} \mathbf{\Sigma} \mathbf{V}^\top, \quad \mathbf{\Sigma} = \text{diag}(\sigma_1, \sigma_2, \dots, \sigma_{r_{\text{max}}})
\end{equation}

توزیع احتمالاتی انرژی مقادیر منفرد به شکل توزیع نرمال‌شده توان مقادیر منفرد تعریف می‌گردد:
\begin{equation}
p_i = \frac{\sigma_i^2}{\sum_{j=1}^{r_{\text{max}}} \sigma_j^2 + \epsilon}, \quad i = 1, \dots, r_{\text{max}}
\end{equation}

آنتروپی شانون طیفی مقادیر منفرد لایه بر اساس توزیع $p$ محاسبه می‌شود:
\begin{equation}
H(\mathbf{\Sigma}) = - \sum_{i=1}^{r_{\text{max}}} p_i \ln(p_i)
\end{equation}

آنتروپی نرمال‌شده $\bar{H} \in [0, 1]$ از طریق مقایسه با بیشینه آنتروپی تئوریک $\ln(r_{\text{max}})$ استخراج می‌گردد:
\begin{equation}
\bar{H} = \frac{H(\mathbf{\Sigma})}{\max(\ln(r_{\text{max}}), \epsilon)}
\end{equation}

رتبه تخصیص‌یافته به لایه بر پایه آنتروپی محاسبه شده و با اعمال کف رتبه مجاز $r_{\text{min}}$ مقید می‌شود:
\begin{equation}
r = \max\left(r_{\text{min}}, \min\left(r_{\text{max}}, \lfloor \bar{H} \cdot r_{\text{max}} \rceil\right)\right)
\end{equation}
در صورت تنظیم رتبه سقف $r_{\text{target}}$ در تنظیمات، مقدار $r \le r_{\text{target}}$ اعمال خواهد شد. فاکتورهای خطی با توزیع متقارن جذر مقادیر منفرد محاسبه می‌شوند:
\begin{equation}
\mathbf{U}_r = \mathbf{U}_{:, :r} \mathbf{\Sigma}_{:r}^{1/2} \in \mathbb{R}^{d_{\text{out}} \times r}, \quad \mathbf{V}_r = \mathbf{\Sigma}_{:r}^{1/2} \mathbf{V}_{:, :r}^\top \in \mathbb{R}^{r \times d_{\text{in}}}
\end{equation}
و لایه متراکم با ماژول دو‌مرحله‌ای \lr{SVDLinear} جایگزین می‌گردد.

\subsection{معماری و پیاده‌سازی}
کلاس \lr{MGIECompressor} در فایل \lr{Codes/framework/compression\_engine/mgie\_svd.py} پیاده‌سازی شده و تحت شناسه‌های \lr{mgie}، \lr{mgie\_svd}، \lr{mgie\_factory} و \lr{mgie\_svd\_factory} در رجیستری فریم‌ورک ثبت شده است. نکات محوری پیاده‌سازی عبارتند از:
\begin{itemize}
    \item \textbf{محاسبه محلی و مستقل آنتروپی (\lr{Per-Layer Entropy Computation}):} بر خلاف روش‌های سرتاسری، در این ماژول هر لایه بر اساس توزیع مقادیر منفرد درونی خود رتبه‌بندی می‌شود که این امر فرآیند فشرده‌سازی را بدون نیاز به حل مسایل بهینه‌سازی سراسری و با حداقل مصرف حافظه انجام می‌دهد.
    \item \textbf{جایگزینی مستقیم با لایه فاکتورگیری‌شده فیزیکی \lr{SVDLinear}:} لایه‌های خطی و لایه‌های \lr{Conv1D} جای خود را به ساختار ماژولار \lr{SVDLinear} می‌دهند که کاهش واقعی پارامترها و تسریع در ضرب ماتریسی را به دنبال دارد.
    \item \textbf{سازگاری با توجه نامتقارن (\lr{GQA Asymmetry Protection}):} با استفاده از تابع \lr{\_is\_gqa\_asymmetric}، در صورت تفاوت در تعداد سرهای کلید/مقدار نسبت به پرسمان، لایه‌های کلید و مقدار از فشرده‌سازی آسیب‌زننده محافظت می‌شوند.
    \item \textbf{آستانه خطای فروبنیوس و مکانیزم برگشت ایمن (\lr{Fallback Guard}):} در صورتی که خطای نسبی فروبنیوس تقریب لایه از آستانه مجاز ($\text{max\_error}=0.50$) فراتر رود، وزن اولیه حفظ شده و لایه جایگزین نمی‌گردد.
\end{itemize}

\subsection{نتایج تجربی}
نتایج ارزیابی رسمی ماژول \lr{mgie\_svd} بر روی مدل‌های بنچ‌مارک بر اساس داده‌های تاییدشده در \lr{thesis\_final/master\_verified\_results.json} در جدول \ref{tab:benchmark_mgie_svd} درج شده است:

\begin{table}[H]
\centering
\small
\caption{نتایج تجربی و اعتبارسنجی بنچ‌مارک فشرده‌ساز \lr{mgie\_svd} روی مدل‌های پایه}
\label{tab:benchmark_mgie_svd}
\begin{tabular}{lccccc}
\toprule
\textbf{مدل پایه} \lr{(Base Model)} & \textbf{دقت} \lr{(Accuracy)} & \textbf{شاخص} \lr{(F1-Score)} & \textbf{حافظه} \lr{(Peak VRAM)} & \textbf{توان} \lr{(Throughput)} & \textbf{وضعیت} \\
\midrule
\lr{DistilBERT} & \lr{90.48\%} & \lr{0.9077} & \lr{311.19 MB} & \lr{182,179.02 tok/s} & \lr{PASS} \\
\lr{RoBERTa} & \lr{72.02\%} & \lr{0.7724} & \lr{301.72 MB} & \lr{107,325.79 tok/s} & \lr{PASS} \\
\lr{GPT-2} & \lr{60.09\%} & \lr{0.6463} & \lr{637.29 MB} & \lr{70,762.65 tok/s} & \lr{PASS} \\
\lr{Qwen} & \lr{---} & \lr{---} & \lr{---} & \lr{---} & \lr{SKIP} \\
\bottomrule
\end{tabular}
\end{table}

داده‌های تجربی گویای دقت ممتاز ماژول بر روی مدل \lr{DistilBERT} با ثبت دقت \lr{90.48\%} و شاخص F1 معادل \lr{0.9077} در کنار سرعت گذر بالای \lr{182,179.02 tok/s} است. در مدل‌های \lr{RoBERTa} و \lr{GPT-2} به ترتیب دقت‌های \lr{72.02\%} و \lr{60.09\%} با وضعیت قبولی (\lr{PASS}) حاصل شده است. ارزیابی مدل \lr{Qwen} در داده‌های ثبت‌شده این شاخه اجرا نشده (\lr{SKIP}) و به عنوان عدم پوشش در بنچ‌مارک علامت‌گذاری شده است.

\subsection{محدودیت‌ها و تحلیل شکست}
تحلیل ساختاری ماژول حاکی از نکات و محدودیت‌های زیر است:
\begin{itemize}
    \item \textbf{فقدان بهینه‌سازی سراسری بودجه رتبه:} تخصیص رتبه به صورت محلی و مستقل بر اساس آنتروپی هر لایه انجام می‌گیرد. این امر می‌تواند منجر به مصرف نامتوازن بودجه پارامتری کل مدل گردد، چرا که وزن لایه‌ها و میزان تاثیرگذاری لایه‌های ابتدایی در مقایسه با لایه‌های انتهایی شبکه در نظر گرفته نمی‌شود.
    \item \textbf{وابستگی دقت دیکودر به حفظ مقادیر منفرد دم طیف:} در شبکه‌های تولیدی نظیر \lr{GPT-2}، حذف مقادیر منفرد با آنتروپی پایین ممکن است به تخریب برخی قابلیت‌های ریززبانی و کاهش شاخص F1 به \lr{0.6463} بیانجامد.
    \item \textbf{افت کارایی در ماتریس‌های با آنتروپی بالا در تمام لایه‌ها:} در صورتی که توزیع مقادیر منفرد تمامی لایه‌ها پهن و با آنتروپی نزدیک به بیشینه باشد، ماژول رتبه‌های بسیار بالایی را به تمامی لایه‌ها اختصاص داده و میزان فشرده‌سازی مورد نظر محقق نخواهد شد.
\end{itemize}
